\documentclass[12pt]{article}
\usepackage{amsmath, amssymb, enumitem, fancyhdr, graphicx, libertine, nth,
  pgfplots, tikz, xcolor}
\pgfplotsset{compat=1.11}
\usetikzlibrary{decorations.pathreplacing}
\usetikzlibrary{patterns}
\usepackage[margin=1in]{geometry}
\usepackage[libertine]{newtxmath}
\pagestyle{fancy}
\definecolor{solution-color}{HTML}{000080}
\chead{\emph{EC201 -- Problem Set 5}}

% Define new topic command
\newcommand{\topic}[1]{\medskip\begin{center}\large\textsc{#1}\normalsize
  \end{center}}

% Define question counter:
\newcounter{question}[section]
\newcommand{\question}[1]{\refstepcounter{question}
  \subsubsection*{Question~\thequestion~-- #1}}

% Subquestions are enumerate with letters
\newenvironment{subquestions}{
  \begin{enumerate}[label=(\roman*)]}{\end{enumerate}}

\begin{document}

\topic{Price Discrimination}
\question{Third-Degree Price Discrimination}
There is one movie theater in a town. The town has two groups of people:
students and non-students. The non-students have a demand curve:
\[
  D_1\left( p_1\right) = 10-p_1
\]
and the students have a demand curve:
\[
  D_2\left( p_2 \right) = 10 - 2p_2
\]
The movie theater has a marginal cost of \$2 per customer (cleaning up their
spilled soda and popcorn). The fixed cost is 10.
\begin{subquestions}
  \item Suppose first the monopolist movie theater cannot price discriminate.
    What price should they charge and what quantity will it sell? What
    profit does it make?
  \item Suppose now the monopolist is able to price discriminate. What price
    should it charge students and what price should it charge non-students?
    How many tickets will they sell?
\end{subquestions}
\question{Price Discrimination}
A monopolist airline faces two types of consumers who value flight quality
differently. There is a high willingness-to-pay group (business people) and a
low willingness-to-pay group (vacationers). The airline faces zero marginal
cost (putting one more person on a plane doens't any extra).

``Quality" is measured on a 30-point scale, where 30 is extremely
luxurious and 0 is extremely uncomfortable.

The two solid downward-sloping lines below measure the willingness-to-pay for a
marginal increase in quality for both groups.  For any quality level, $q$, the
willingness-to-pay is the area \emph{under the line} between between 0 and $q$. Each triangle blow has an area of 50 and each square has an area of 100.
If $q=10$, vacationers are willing to pay at most $150$ for air travel,
while business people would pay $250$.  If the quality level is 0,
the area is always 0 so no one would pay for airfare at quality 0.

\begin{center}
  \begin{tikzpicture}[scale=0.8]
    \draw[thick,<->] (0,9) node[above] {$WTP$}--(0,0)--(9,0) node[below right]
    {Quality, $q$};
    \draw[dashed,-] (0,5)--(2.5,5);
    \draw[dashed,-] (5,0)--(5,2.5);
    \draw[dashed,-] (2.5,0)--(2.5,5);
    \draw[dashed,-] (0,2.5)--(5,2.5);
    \draw (0, 7.5) node[left] {30};
    \draw (0, 5) node[left] {20};
    \draw (0, 2.5) node[left] {10};
    \draw (2.5, 0) node[below] {10};
    \draw (5, 0) node[below] {20};
    \draw (7.5, 0) node[below] {30};
    \draw (0, 0) node[below left] {0};
    \draw (0,7.5)--(7.5,0) node[above right] {\textsl{Business people}};
    \draw (0,5)--(5,0) node[above right] {\textsl{Vacationers}};
    % \draw (0.8, 6.2) node {$A$};
    % \draw (1, 3.25) node {$B$};
    % \draw (1.5, 4.25) node {$C$};
    % \draw (3.2,3.75) node {$D$};
    % \draw (1,1) node {$E$};
    % \draw (3,1) node {$F$};
    % \draw (4,2) node {$G$};
    % \draw (6,1) node {$H$};
  \end{tikzpicture}

\end{center}
\begin{subquestions}
  \item If the only people in the world were business people, what should the
    airline charge for air travel to maximize profits?
  \item If the only people in the world were vacationers, what should the
    airline charge for air travel to maximize profits?
\end{subquestions}
For the rest of the question, assume that 50\% of air travelers are business
people and 50\% are vacationers.
\begin{subquestions}
  \setcounter{enumi}{2}
  \item Suppose vacationers were able to obtain a ``vacationer's card'' which
    allowed them to get discounts for air travel if discounts were available
    (like a student card). The airline now has an observable way to tell
    travelers apart. What should the airline charge each group to maximize
    profits? What is the average profit made per person?\footnote{If you get
      revenue $x$ from vactioners and revenue $y$ from business people, the
      average profit per person is $\frac{x+y}{2}$. If you only sell to
      business people, the average profit per person is $\frac{y}{2}$.}
  \item Suppose now that a vacationer's card doesn't exist so the airline can't
    charge different people different prices.
    If the airline were to choose only one price-quality plan, what would
    it do to maximize profits? To do this you should check two possibilities:
    \begin{itemize}
      \item You set a price-quality plan such that both groups would be
        willing to buy.
      \item You set a price-quality plan such that only business people will
        be willing to buy (but you lose out on half the market).
    \end{itemize}
    Find the average profit made per person from both of these possibilities.
  \item Suppose again that the vacationer's card doesn't exist but the airline
    now wants to set two price-quality plans in order to make the vacationers
    and business people \emph{self-select} into each plan.
    \begin{itemize}
      \item What price-quality plans for both groups would make the groups
        self-select and maximize profits?
      \item What is the average profit per customer?
    \end{itemize}
\end{subquestions}
% \question{Bundling}
% You are a monopolist software company sell three different software products.
% Three types of people buy your software: Marketers, Accountants and Graphic
% Designers. There is an equal amount of each type of person in the population.
% The willingness to pay of each group for each type software is below:
% \begin{center}
%   \begin{tabular}{|l|c|c|c|}\hline
%                            & Marketers   & Accountants & Graphic Designers \\
%     \hline
%     Presentation Software  & 150         & 100         & 0                 \\
%     Spreadsheet Software   & 100         & 150         & 0                 \\
%     Image Editing Software & 40          & 10          & 150               \\
%     \hline
%   \end{tabular}
% \end{center}
% You are not able to price discriminate between the different groups.
% \begin{subquestions}
%   \item If you were selling each piece of software individually, what price
%     should you charge for each?
%   \item If you were able to bundle two or three pieces of software together,
%     what pieces of software would you bundle together?
% \end{subquestions}
\question{Two-Part Tariffs}
There is a single theme park in a large state. It can put people on a roller
coaster at zero marginal cost (but they do have a fixed cost). The demand for
each person for roller coaster rides is given by $D\left( p \right)=200-\frac{1}{2}p$, where
$p$ is the price of each roller coaster ride. What price should
the theme park charge for entry into the theme park and how much should it
charge for each person to ride the roller coaster?
\topic{Game Theory}
\question{Simulataneous Games}
Find all Nash equilibria (pure and mixed) of the following games:
\begin{subquestions}
  \item
  \setlength{\arraycolsep}{0pt}
  \[
  \begin{array}{cccccc}
  &  & \quad L \quad & & \quad R \quad  &  \\
  \cline{3-6} U \quad & \vline{} & \quad 2,4  \quad & \vline{} & \quad 6,3 \quad & \vline{}  \\
  \cline{3-6} D \quad & \vline{} & \quad 3,1 \quad & \vline{} & \quad 5,3 \quad & \vline{} \\
  \cline{3-6}
  \end{array}
  \]
  \item
  \setlength{\arraycolsep}{0pt}
  \[
  \begin{array}{cccccc}
  &  & \quad L \quad & & \quad R \quad  &  \\
  \cline{3-6} U \quad & \vline{} & \quad 2,1  \quad & \vline{} & \quad 0,0 \quad & \vline{}  \\
  \cline{3-6} D \quad & \vline{} & \quad 0,0 \quad & \vline{} & \quad 1,2 \quad & \vline{} \\
  \cline{3-6}
  \end{array}
  \]
  \item
  \setlength{\arraycolsep}{0pt}
  \[
  \begin{array}{cccccc}
  &  & \quad L \quad & & \quad R \quad  &  \\
  \cline{3-6} U \quad & \vline{} & \quad -2,-2  \quad & \vline{} & \quad 4,0 \quad & \vline{}  \\
  \cline{3-6} D \quad & \vline{} & \quad 0,4 \quad & \vline{} & \quad 2,2 \quad & \vline{} \\
  \cline{3-6}
  \end{array}
  \]
\end{subquestions}
\question{Sequential Games}
Consider the following extensive form game:
\ \\
\begin{center}
\begin{tikzpicture}[scale=1.5,font=\footnotesize]
\tikzstyle{solid node}=[circle,draw,inner sep=1.5,fill=black]
\tikzstyle{hollow node}=[circle,draw,inner sep=1.5]
\tikzstyle{level 1}=[level distance=15mm,sibling distance=3.5cm]
\tikzstyle{level 2}=[level distance=15mm,sibling distance=1.5cm]
\tikzstyle{level 3}=[level distance=15mm,sibling distance=1cm]
\node(0)[hollow node,label=above:{$P1$}]{}
child{node[solid node,label=above left:{$P2$}]{}
child{node[label=below:{$(5,3)$}]{} edge from parent node[left]{$A$}}
child{node[label=below:{$(0,2)$}]{} edge from parent node[left]{$B$}}
edge from parent node[left,xshift=-5]{$L$}
}
child{node[label=below:{$(4,2)$}]{} edge from parent node [above]{$R$}
};
\end{tikzpicture}
\end{center}
\begin{subquestions}
  \item Write the normal form (payoff matrix) of the game.
  \item Find all the Nash equilibria of the game.
  \item Which of the Nash equilibria that you found are subgame perfect?
\end{subquestions}
\topic{Oligopoly}
\question{Different Market Structures}
The demand curve for a good is given by $D\left( p
\right)=5-\frac{1}{2}p$. The cost function for a firm is given by $c\left( q
\right)=2q$.
\begin{subquestions}
  \item The market was served by perfectly competitive firms.
      What is the market price and total output produced by the industry?  What
      profit does each firm make?
  \item The market is served by a monopolist. What price and quantity does the
    monopolist produce? What profit does it make?
  \item The market is served by two firms who compete simultaneously on price.
    What price does each firm charge? What quantity is produced by both firms?
    What profit does each firm make?
  \item  The market is served by two firms who simulatenously choose their
    quantites. The price is then set to clear the market. What quantity is
    produced by both firms? What is the price? What is the profit for each
    firm?
  \item The market is served by two firms but one firm chooses the quantity to
    produce first. The other firm then observes this quantity and chooses its
    own quantity. The price is then set to clear the market. What quantity does
    each firm produce? What is the price? What profit does each firm make?
\end{subquestions}


\end{document}
